\section{Method}
\label{sec:method}

We describe the final system top-down, as it would be reproduced, rather than in the order
we discovered it; the chronological development trace of every direction we explored is given
in Appendix~\ref{app:experiment_summary}. The system has three parts. First, a training corpus
of single-plant images carrying taxonomic labels (\S\ref{sec:cap_extra}). Second, a partial
fine-tune of BioCLIP~2.5 ViT-H/14 that adapts only the last four transformer blocks and
attaches auxiliary taxonomy heads, preserving the Tree-of-Life prior that makes the backbone
valuable on biological data (\S\ref{sec:backbone}--\S\ref{sec:bioclip25_010}). Third, a
test-time pipeline that tiles each quadrat, pools the per-tile votes by softmax mean,
calibrates with class-prior logit adjustment, and selects an adaptive number of species
(\S\ref{sec:adaptive_inference}--\S\ref{sec:calibration}). We close with a single canonical
statement of the final configuration (\S\ref{sec:final_submission}). Throughout this
section we define \emph{what} the system is and defer all leaderboard numbers to
Section~\ref{sec:experiments}.

\subsection{Data and Label Construction}
\label{sec:cap_extra}

The training data are single-plant images from the PlantCLEF~2024 corpus, each contributed
through the Pl@ntNet citizen-science pipeline and carrying one species label together with its
Linnaean ancestry (genus, family, and, where available, order and class). The distribution is
heavily long-tailed: a small number of common species exceed one thousand images while hundreds
of rare species carry only a handful, so vanilla shuffled training exposes the network to
roughly two orders of magnitude more head-class than tail-class examples per epoch.

Our final model trains on an enlarged and cleaned manifest. The original PlantCLEF~2024 metadata
(merged with a GBIF taxonomy lookup) supplies the first fine-tune; a subsequent data rebuild
(experiment~i001) re-downloaded the corpus, supplemented it with a research-grade iNaturalist
pull, deduplicated, and pre-filled genus and family labels for every row, yielding roughly
$2.65$\,million images across $7{,}806$ species (Table~\ref{tab:i001_manifest}). Our final
model (experiment~i002) trains on this larger manifest \emph{without} any per-species cap. We
also explored an explicitly rebalanced variant (experiment~i003) that augments the tail with
extra validated images and caps every species at five hundred images; its leaderboard outcome,
which did not improve on the uncapped manifest, is analysed in Section~\ref{sec:longtail}. We
therefore treat ``more clean data, no cap'' as a deliberate design choice for the final system
rather than as a controlled study of rebalancing.

\begin{table}[ht]
  \centering
  \small
  \setlength{\tabcolsep}{6pt}
  \begin{tabular}{@{}lrrrr@{}}
    \toprule
    Source / Split & Images & Observations & Species & Genera \\
    \midrule
    PlantCLEF~2024 (Pl@ntNet)     & $1{,}408{,}033$ & $1{,}151{,}904$ & $7{,}806$ & $1{,}446$ \\
    iNaturalist (research-grade)  & $1{,}245{,}748$ & $1{,}245{,}085$ & $5{,}037$ & $1{,}254$ \\
    \midrule
    \textbf{i001 manifest (union)}    & $\mathbf{2{,}653{,}781}$ & --- & $\mathbf{7{,}806}$ & $\mathbf{1{,}446}$ \\
    \hspace{1em}Train (stratified 90\%) & $2{,}391{,}457$ & $1{,}121{,}447$ & $7{,}806$ & $1{,}446$ \\
    \hspace{1em}Val (stratified 10\%)   & $262{,}324$     & $123{,}768$     & $7{,}309$ & $1{,}429$ \\
    \bottomrule
  \end{tabular}
  \caption{Composition of the i001 training manifest used by the i002 model. The PlantCLEF~2024 row reproduces the dataset overview from~\cite{joly2024overview}; the iNaturalist row reflects the research-grade pull we de-duplicated and image-verified before adding to the manifest. All iNaturalist species are a subset of the PlantCLEF~2024 taxonomy, so the union species and genus counts match the PlantCLEF~2024 row; iNaturalist therefore covers $5{,}037$ of the $7{,}806$ classes (the most populous head) rather than extending the long tail. The Train / Val rows are the deterministic stratified split applied at training time (10\% per species, seed $42$, species with fewer than five images held entirely in Train); the Val row drops to $7{,}309$ species because $497$ rare species fall below the five-image threshold and remain Train-only. Observations are reported only on the union and split rows where the field is well defined for iNaturalist sources, and they are omitted on the union total because the field is undefined for PlantCLEF~2024 rows in our merged manifest.}
  \label{tab:i001_manifest}
\end{table}

\subsection{BioCLIP 2.5 Fine-tuning}
\label{sec:backbone}

BioCLIP~\cite{stevens2024bioclip} is a contrastive vision-language model whose pretraining on
TreeOfLife-10M aligns its feature space with the Linnaean taxonomy, which is precisely the prior
that makes it valuable for biological classification. We adapt the BioCLIP~2.5 ViT-H/14 image
encoder by supervised fine-tuning, but only partially: rather than updating the full backbone,
we freeze the lower twenty-eight transformer blocks and unfreeze only the last four blocks
together with the final layer norm and projection layer. This narrow trainable surface is enough
to adapt the encoder to the single-plant distribution while preserving the broader biological
prior, and it is more resilient to optimiser noise on a long-tailed training set than full
fine-tuning. The number of unfrozen blocks is itself a sharp design variable; we report the
ablation that locates its optimum in Section~\ref{sec:experiments}.

Training is staged in two phases. An initial ten epochs of head-only training on the enlarged
manifest with the backbone fully frozen initialises the per-head MLPs of
\S\ref{sec:bioclip25_010}; a further ten epochs with the last four transformer blocks unfrozen
then produce the final i002 checkpoint, resuming the model weights only. Both stages train in
bfloat16 across two RTX~5090 GPUs under cross entropy with linear warmup and cosine decay; the
complete hyperparameter specification (learning rates, batch sizes, gradient accumulation,
augmentations, image transforms, and the validation split) is given in
Appendix~\ref{sec:training_config}, Table~\ref{tab:training_config}.

\subsection{Taxonomy-aware Heads}
\label{sec:bioclip25_010}

On top of the encoder we attach a small head that predicts species jointly with coarser
taxonomic levels, so the auxiliary taxonomy losses act as a structural regulariser. Each head is
built from the same template (layer normalisation of the $1{,}024$-dimensional CLS embedding, a
linear projection to a $1{,}024$-dimensional hidden representation, a GELU non-linearity, and
dropout at rate $0.2$), but the two fine-tunes we report wire it differently.

\textbf{Shared trunk (experiment~010).} Our first strong fine-tune uses a single \emph{shared}
multi-layer perceptron whose one hidden vector is read by every taxonomic head, so the auxiliary
losses regularise the exact representation the species head consumes. It exposes five heads, with
output dimensions $7{,}806$ (species), $1{,}446$ (genus), $181$ (family), $61$ (order), and $6$
(class).

\textbf{Per-head MLPs (experiment~i002).} Experiment~i002 instead gives each head its own
\emph{independent} multi-layer perceptron of the same structure, decoupling the per-task
gradients above the backbone. Because the redownloaded i001 manifest carries species, genus, and
family columns but no order or class columns, it retains three per-head branches and drops the
two coarsest levels.

The training objective is a weighted joint cross entropy
\begin{equation}
  \mathcal{L} \;=\; \mathcal{L}_{\text{species}} \;+\; \!\!\!\!\sum_{l \in L}\!\!\!\! w_l\,\mathcal{L}_l,
\end{equation}
where the species term carries an implicit weight of one, $L = \{g,f,o,c\}$ for the shared-trunk
fine-tune and $L = \{g,f\}$ for the i002 per-head model, and the auxiliary weights are
$w_g{=}0.30$, $w_f{=}0.15$, $w_o{=}0.05$, $w_c{=}0.02$. These are fixed, hand-set constants
rather than learned or entropy-derived weights; coarser levels are deliberately down-weighted so
that, when present, the order and class heads act primarily as regularisation rather than driving
the encoder. The i002 per-head model uses only $w_g{=}0.30$ and $w_f{=}0.15$. Missing taxonomy
labels are encoded as $-1$ and masked out of the auxiliary loss.

We stress that the shared-trunk and per-head designs differ in several other respects at the same
time (training data, schedule, number of unfrozen blocks, and head count), so we do \emph{not}
run a controlled comparison that varies only the head wiring, and we make no claim that the
per-head MLP, or the taxonomy supervision, caused any performance gain. We report the head wiring
as a deliberate design choice and analyse the resulting feature spaces descriptively in
Section~\ref{sec:head_repr}.

\begin{figure*}[t]
  \centering
  \resizebox{0.95\linewidth}{!}{% Head-design schematic -- 010 shared MLP trunk vs i002 independent per-head MLPs.
% Pictorial only; dims from report/head_features/tex_info (hidden 1024; heads
% species 7806, genus 1446, family 181, order 61, class 6).
% Style follows figures/aggregation_definitions_figure.tex (sffamily, accent
% colours, rounded cards). Include in a figure* float (see sections/experiments.tex).
%
% Required preamble (already in report/preamble.tex):
%   \usepackage{tikz}
%   \usetikzlibrary{positioning,arrows.meta,calc,fit,backgrounds}

\begin{tikzpicture}[
  font=\sffamily,
  >={Stealth[length=4.5pt,width=5pt]},
  card/.style={rounded corners=4pt},
  enc/.style={draw=Enc!55, fill=Enc!12, rounded corners=3pt, line width=0.6pt,
              align=center, font=\sffamily\scriptsize, inner sep=4pt},
  mlp/.style={draw=#1!60, fill=#1!14, rounded corners=3pt, line width=0.6pt,
              align=center, font=\sffamily\scriptsize, inner sep=3.5pt, minimum height=8mm},
  head/.style={draw=#1!55, fill=#1!8, rounded corners=2pt, line width=0.5pt,
               align=center, font=\sffamily\scriptsize, inner sep=2.4pt, minimum height=7mm},
  flow/.style={->, line width=0.6pt, draw=black!55},
  grad/.style={->, densely dashed, line width=0.6pt, draw=#1!75},
  collab/.style={font=\sffamily\small\bfseries, text=black!70},
  note/.style={font=\sffamily\scriptsize, text=black!60, align=center},
]

% --- palette -----------------------------------------------------------------
\definecolor{Enc}{RGB}{70,90,120}      % backbone slate
\definecolor{Shar}{RGB}{123,84,193}    % 010 shared trunk (violet)
\definecolor{Sp}{RGB}{33,102,172}      % species (blue)
\definecolor{Ge}{RGB}{42,157,123}      % genus (teal-green)
\definecolor{Fa}{RGB}{226,135,38}      % family (amber)
\definecolor{Oc}{RGB}{120,120,120}     % order/class (grey)

% =============================================================================
% PANEL A -- experiment 010: shared MLP trunk -> five heads
% =============================================================================
\begin{scope}[shift={(0,0)}]
  \node[collab] at (3.2,4.55) {Experiment 010: shared MLP trunk};

  % backbone
  \node[enc, minimum width=6.2cm, minimum height=8mm] (b010) at (3.2,0)
    {BioCLIP 2.5 ViT-H/14 encoder};

  % shared MLP
  \node[mlp=Shar, minimum width=4.2cm] (s010) at (3.2,1.55)
    {Shared MLP};

  % five heads
  \node[head=Sp] (h010sp) at (0.55,3.35) {species\\\tiny 7806};
  \node[head=Ge] (h010ge) at (1.85,3.35) {genus\\\tiny 1446};
  \node[head=Fa] (h010fa) at (3.15,3.35) {family\\\tiny 181};
  \node[head=Oc] (h010or) at (4.45,3.35) {order\\\tiny 61};
  \node[head=Oc] (h010cl) at (5.75,3.35) {class\\\tiny 6};

  % flow
  \draw[flow] (b010) -- (s010);
  \draw[flow] (s010.north) -- (h010sp.south);
  \draw[flow] (s010.north) -- (h010ge.south);
  \draw[flow] (s010.north) -- (h010fa.south);
  \draw[flow] (s010.north) -- (h010or.south);
  \draw[flow] (s010.north) -- (h010cl.south);

  % gradient coupling: taxonomy losses pass back through the shared trunk
  \draw[grad=Shar] (h010cl.east) to[out=0,in=0] (s010.east);
  \node[note, text=Shar!75!black, text width=6.2cm, anchor=center] at (3.2,4.05)
    {dashed: taxonomy-loss gradients flow back through the shared trunk};
\end{scope}

% =============================================================================
% PANEL B -- experiment i002: independent per-head MLPs
% =============================================================================
\begin{scope}[shift={(8.0,0)}]
  \node[collab] at (2.6,4.55) {Experiment i002: independent per-head MLPs};

  % backbone
  \node[enc, minimum width=6.2cm, minimum height=8mm] (bi02) at (2.6,0)
    {BioCLIP 2.5 ViT-H/14 encoder};

  % three per-head MLPs
  \node[mlp=Sp, minimum width=1.5cm] (mi02sp) at (0.8,1.55) {species\\MLP};
  \node[mlp=Ge, minimum width=1.5cm] (mi02ge) at (2.6,1.55) {genus\\MLP};
  \node[mlp=Fa, minimum width=1.5cm] (mi02fa) at (4.4,1.55) {family\\MLP};

  % three heads
  \node[head=Sp] (hi02sp) at (0.8,3.35) {species\\\tiny 7806};
  \node[head=Ge] (hi02ge) at (2.6,3.35) {genus\\\tiny 1446};
  \node[head=Fa] (hi02fa) at (4.4,3.35) {family\\\tiny 181};

  % flow: backbone -> each MLP -> own head (gradients stay in branch)
  \draw[flow] (bi02.north -| mi02sp) -- (mi02sp.south);
  \draw[flow] (bi02.north -| mi02ge) -- (mi02ge.south);
  \draw[flow] (bi02.north -| mi02fa) -- (mi02fa.south);
  \draw[flow] (mi02sp) -- (hi02sp);
  \draw[flow] (mi02ge) -- (hi02ge);
  \draw[flow] (mi02fa) -- (hi02fa);

  \node[note, text=black!60, text width=5.0cm, anchor=north] at (2.6,-0.7)
    {auxiliary gradients stay within each branch above the shared backbone};
\end{scope}

\end{tikzpicture}
}
  \caption{Head designs of the two supervised fine tunes. Experiment~010 (left) uses a single shared MLP (layer normalisation, a linear projection of the $1{,}024$-d embedding, GELU, and dropout) whose hidden vector feeds five linear heads (species/genus/family/order/class); the auxiliary taxonomy losses back propagate through the shared trunk. Experiment~i002 (right) uses three independent per head MLPs (species/genus/family), so the auxiliary gradients do not pass through the species MLP. Both share the partially unfrozen BioCLIP~2.5 backbone (last blocks plus the final layer norm and projection). The diagram shows the design difference only: the two experiments are not a controlled ablation, since they also differ in training data, schedule, unfrozen blocks, and head count.}
  \label{fig:head_design}
\end{figure*}

\subsection{Tiled Inference}
\label{sec:adaptive_inference}

The classifier is trained on single-plant images, but each test quadrat is a dense, multi-species
scene in which plants overlap, occlude one another, and appear at varying scales. Passing a whole
quadrat through an encoder trained on isolated specimens therefore mixes many species into one
embedding. We instead partition each quadrat into a non-overlapping
$4\times4$ grid of sixteen tiles, each $1/4$ of the quadrat in width and height (and $1/16$ of
its area), so that most tiles contain only one or a few species, closer to the training
distribution. Each tile is forwarded through the encoder
at two resolutions, the encoder's native $224$ pixels and $336$ pixels (the ViT-H/14 backbone
tolerates the $1.5\times$ stretch via bicubic positional-embedding interpolation), giving a
per-tile species distribution that the next stage pools.

\subsection{Tile Aggregation and Selection}
\label{sec:tile_agg_method}

The sixteen per-tile predictions must be pooled into one image-level score vector before we
select the species to emit. We implemented three pooling rules. \emph{Mean} averages the per-tile
logits; it is robust to a single noisy tile but dilutes a species that is strong in only one or
two tiles among the many background tiles where it is absent. \emph{Max} keeps each species'
single strongest tile; it suits a species confined to one tile but is fragile, since one
over-confident background tile can promote a species across the whole image. \emph{Softmax-mean}
turns each tile into a normalised probability vote and averages those votes, so no single tile can
dominate while evidence still accumulates across tiles. The final system uses softmax-mean. Plain mean was explored on an earlier head-only checkpoint
(experiment~010) and dropped; it was not re-tested on the final model. The two-stage comparison
behind this choice, with the same-checkpoint caveat, is reported in
Section~\ref{sec:tile_aggregation}.

After aggregation we replace standard top-$k$ selection with an adaptive thresholded rule: we emit
every species whose probability exceeds a fixed threshold $T$, subject to a floor of $k_{\min}{=}2$
and a ceiling of $k_{\max}{=}10$ species per quadrat. This matches the test set's empirical
cardinality of roughly two to three species per quadrat without committing to a single $k$. The
operating point ($T{=}0.03$) and its comparison against fixed top-$k$, gap, and relative-threshold
rules are reported in Section~\ref{sec:experiments}.

\subsection{Calibration and Ensembling}
\label{sec:calibration}

Two further ingredients complete the inference recipe. First, to counteract the long-tail bias of
the species prior we apply class-prior logit adjustment~\cite{menon2021longtail},
$\ell_s' = \ell_s - \tau \log \tilde{\pi}_s$ with $\tau{=}0.25$, where $\tilde{\pi}_s$ is the
Laplace-smoothed training prior for species $s$; this requires larger margins for common species
and smaller margins for rare ones. Second, the per-tile distributions from the $224$- and
$336$-pixel passes are averaged in probability space before selection, a single-checkpoint
resolution ensemble that recovers fine-scale evidence the native resolution drops. Both knobs are
swept on the leaderboard in Section~\ref{sec:experiments}; here we fix the values used by the
final system.

\subsection{Final System}
\label{sec:final_submission}

Our final system is a single BioCLIP~2.5 ViT-H/14 backbone trained under the two-stage
schedule of \S\ref{sec:backbone} on the enlarged i001 manifest. Only the last four transformer
blocks (plus the final layer norm and projection) are unfrozen in the second stage, and both
stages share the joint species/genus/family cross-entropy loss of \S\ref{sec:bioclip25_010}. At
inference (Figure~\ref{fig:inference_pipeline}), each quadrat is
partitioned into a non-overlapping $4\times4$ grid of sixteen tiles; every tile is forwarded
through the encoder at both $224$ and $336$ pixels. Per-tile softmax probabilities are aggregated
by species-wise mean across the sixteen tiles, the two resolutions are averaged, class-prior logit
adjustment ($\tau{=}0.25$) is applied, and every species whose post-adjusted probability exceeds
$T{=}0.03$ is emitted, subject to $k_{\min}{=}2$ and $k_{\max}{=}10$. The two-resolution average is a single-checkpoint ensemble of the i002 model, not a
multi-checkpoint ensemble; it is applied at inference time, not as part of the training-time
pipeline. This configuration is our
headline submission, scoring $0.41826$ public and $0.40283$ private Macro~F1 on the PlantCLEF~2026
leaderboard; the experiments that justify each component follow in Section~\ref{sec:experiments}.

Beyond this anchor we submitted three architecture-orthogonal pivots that aimed to add information
the visual encoder cannot read from a single tile: a triple-backbone classifier retrained on
frozen BioCLIP~2.5, DINOv3, and ConvNeXt-V2-L features; a genus and family co-occurrence rerank;
and a seasonal phenology prior built from GBIF month-of-observation histograms (full method in
Appendix~\ref{sec:phenology}). None improved over the visual anchor; all three are reported in
Section~\ref{sec:experiments}.

\begin{figure*}[t]
  \centering
  \includegraphics[width=0.95\linewidth]{figures/inference_pipeline}
  \caption{Inference pipeline for the ANU PlantCLEF 2026 submission
(public Macro F1 $0.418265$). Each quadrat is partitioned into a
non-overlapping $4{\times}4$ grid of sixteen $448$-pixel tiles and
forwarded through the partially fine-tuned BioCLIP~2.5 ViT-H/14
(lower 28 transformer blocks frozen; last 4 plus LN and projection
unfrozen) at two inference resolutions, $224$ and $336$ pixels.
The encoder produces outputs at three taxonomic levels (family,
genus, species); only the species head is used at inference time.
Per-tile softmax probabilities from the species head are averaged
across the 16 tiles and across the two resolutions, then class-prior
logit adjustment ($\tau{=}0.25$) down-weights common species before
an adaptive probability threshold ($T{=}0.03$, $k_{\min}{=}2$,
$k_{\max}{=}10$) selects the output species set.}
  \label{fig:inference_pipeline}
\end{figure*}
